% ===== PRÉAMBULE CANONIQUE — à coller VERBATIM en tête de CHAQUE .tex =====
\documentclass[11pt,a4paper]{article}
\usepackage[utf8]{inputenc}\usepackage[T1]{fontenc}\usepackage{lmodern}
\usepackage[french]{babel}
\usepackage{amsmath,amssymb,mathtools}\usepackage{icomma}
\usepackage{siunitx}\sisetup{locale=FR}  % \qty{}{}, \si{}, \ang{}
\usepackage{xcolor}\usepackage{colortbl}
\usepackage{tikz}\usepackage{pgfplots}\pgfplotsset{compat=1.18}
\usepackage[siunitx,europeanresistors]{circuitikz} % circuits (élec.)
\usepackage[most]{tcolorbox}\usepackage{enumitem}\usepackage{array,booktabs}
\usepackage[a4paper,margin=2.2cm]{geometry}
\usetikzlibrary{arrows.meta,calc,angles,quotes,positioning,patterns,%
  backgrounds,shapes.geometric,decorations.pathmorphing,decorations.markings,intersections}
\definecolor{primary}{RGB}{222,49,99}\definecolor{primarydark}{RGB}{150,22,55}
\definecolor{defcol}{RGB}{0,114,178}\definecolor{methcol}{RGB}{52,92,148}
\definecolor{excol}{RGB}{0,135,90}\definecolor{attncol}{RGB}{200,40,55}
\tcbset{boxbase/.style={enhanced,breakable,boxrule=0.9pt,arc=2.5pt,
  left=3mm,right=3mm,top=2.3mm,bottom=2.3mm,fonttitle=\bfseries,coltitle=white}}
% --- boîtes TUTORIEL (noms EXACTS, ne pas renommer) ---
\newtcolorbox{cle}[1][]{boxbase,colback=defcol!6,colframe=defcol,
  title={\textsc{À retenir}\ifx\relax#1\relax\relax\else\textmd{ : \itshape #1}\fi}}
\newtcolorbox{methode}[1][]{boxbase,colback=methcol!7,colframe=methcol,sharp corners,
  title={\textsc{Méthode}\ifx\relax#1\relax\relax\else\textmd{ --- \itshape #1}\fi}}
\newtcolorbox{exemplebox}[1][]{boxbase,colback=excol!5,colframe=excol,
  title={\textsc{Exemple}\ifx\relax#1\relax\relax\else\textmd{ : \itshape #1}\fi}}
\newtcolorbox{piege}[1][]{boxbase,colback=attncol!6,colframe=attncol,
  title={\textsc{Piège}\ifx\relax#1\relax\relax\else\textmd{ : \itshape #1}\fi}}
\newtcolorbox{remarque}[1][]{boxbase,colback=black!4,colframe=black!45,
  title={\textsc{Remarque}\ifx\relax#1\relax\relax\else\textmd{ : \itshape #1}\fi}}
% --- boîtes EXERCICE / CORRIGÉ (noms EXACTS, auto-comptées) ---
\newtcolorbox[auto counter]{exo}[1][]{boxbase,colback=primary!4,colframe=primary,
  title={\textsc{Exercice}~\thetcbcounter\ifx\relax#1\relax\relax\else\textmd{ --- \itshape #1}\fi}}
\newtcolorbox[auto counter]{corrige}[1][]{boxbase,colback=excol!4,colframe=excol,
  title={\textsc{Corrigé}~\thetcbcounter\ifx\relax#1\relax\relax\else\textmd{ --- \itshape #1}\fi}}
\newcommand{\vv}[1]{\vec{#1}}\newcommand{\ds}{\displaystyle}
\newcommand{\R}{\mathbb{R}}\newcommand{\N}{\mathbb{N}}\newcommand{\Z}{\mathbb{Z}}
% (un \usepackage{fancyhdr}+en-tête est possible ; il est ignoré côté web)
% =========================================================================
\begin{document}
\sisetup{per-mode=symbol}

\begin{corrige}[coefficient de frottement cinétique]
\begin{enumerate}[label=\alph*)]
  \item Seul le frottement travaille (résistant). À l'arrêt $v_B=0$ :
        \[
        -\mu mg\,d=E_c(B)-E_c(A)=0-\tfrac12 m v^2 .
        \]
  \item La masse se simplifie : $\mu g\,d=\tfrac12 v^2$, donc
        \[
        \mu=\frac{v^2}{2gd}=\frac{4^2}{2\times10\times10}=\frac{16}{200}=0,08 .
        \]
\end{enumerate}
\end{corrige}

\begin{corrige}[comment varie l'énergie cinétique pendant un lancer vertical ?]
\begin{enumerate}[label=\alph*)]
  \item $E_c(t)=\tfrac12 m (v_0-gt)^2=\tfrac12\times0,2\times(10-10t)^2=0,1\,(10-10t)^2$.
        \begin{itemize}[leftmargin=1.4em,topsep=1pt]
          \item à $t=0$ : $E_c=0,1\times10^2=\qty{10}{\joule}$ ;
          \item au sommet $t=\qty{1}{\second}$ : $v=0$, donc $E_c=\qty{0}{\joule}$.
        \end{itemize}
  \item \textbf{Non}, ce n'est pas une droite : $E_c$ dépend de $(v_0-gt)^2$, c'est une \emph{parabole}
        en $t$. Pendant la montée, $E_c$ \emph{décroît} de \qty{10}{\joule} à $0$ ; pendant la descente,
        elle \emph{croît} à nouveau symétriquement. (C'est la vitesse $v(t)$, elle, qui varie
        linéairement — pas l'énergie cinétique.)
\end{enumerate}
\end{corrige}

\begin{corrige}[descente d'un plan incliné avec frottement]
\begin{enumerate}[label=\alph*)]
  \item $h=L\sin\ang{30}=5\times0,5=\qty{2,5}{\metre}$.
        $W(\vv G)=mgh=2\times10\times2,5=\qty{50}{\joule}$ ;
        $W(\vv F_f)=-F_f\,L=-4\times5=\qty{-20}{\joule}$ ; $W(\vv N)=0$.
  \item $\sum W=50-20=\qty{30}{\joule}=\tfrac12 m v_B^2$, donc
        $v_B^2=\dfrac{2\times30}{2}=30$, soit $v_B=\sqrt{30}\approx\qty{5,5}{\metre\per\second}$.
        Sans frottement on aurait $\sqrt{2gh}=\sqrt{2\times10\times2,5}=\sqrt{50}\approx\qty{7,1}{\metre\per\second}$ :
        le frottement a réduit la vitesse d'arrivée.
\end{enumerate}
\end{corrige}

\begin{corrige}[force de freinage moyenne d'une voiture]
\begin{enumerate}[label=\alph*)]
  \item $v=\dfrac{108}{3,6}=\qty{30}{\metre\per\second}$, d'où
        $E_c=\tfrac12\times1000\times30^2=\tfrac12\times1000\times900=\qty{450000}{\joule}=\qty{450}{\kilo\joule}$.
  \item $F=\dfrac{E_c}{d}=\dfrac{450000}{75}=\qty{6000}{\newton}$.
\end{enumerate}
\end{corrige}

\begin{corrige}[vrai ou faux : l'énergie cinétique]
\begin{enumerate}[label=\alph*)]
  \item \textbf{Faux}. $E_c\propto v^2$ : si $v$ double, $E_c$ est multipliée par $2^2=4$ (elle
        quadruple, elle ne double pas).
  \item \textbf{Vrai}. $E_c=\tfrac12 m v^2$ ne dépend pas de $g$ : à $m$ et $v$ identiques, l'énergie
        cinétique est la même partout.
  \item \textbf{Vrai}. $\sum W=\Delta E_c=0$ implique $E_c(B)=E_c(A)$, donc la vitesse (en norme) est
        conservée.
  \item \textbf{Faux}. $E_c=\tfrac12 m v^2\geq0$ toujours (produit d'un carré par une masse positive).
\end{enumerate}
\end{corrige}
\end{document}
